\chapter{Transport, Label Pressure, and Specific VFE}
\label{ch:normalized-geometry}
\label{ch:geometry-pressure}

This chapter asks what becomes of a finite label ensemble of
\cref{def:rg-geometry-ensemble} when every branch is a native lattice law of
\cref{ch:constructive-spatial-rg} and the volume grows: which label survives,
at what rates the evidence and the variational free energy scale, and how
these rates move along a fixed blocking word. The answers are the pressure
with explicit surface rates of \cref{thm:rg-pressure-boundary}, strict label
selection with the exact tie condition of \cref{thm:rg-pressure-selection},
the specific VFE of \cref{thm:rg-pressure-vfe-density} with its minimum,
minus the evidence density, in \cref{prop:rg-pressure-vfe-minimum}, and the
transport of the specific VFE along a fixed word in
\cref{thm:rg-pressure-induced-vfe}. The chapter opens with the finite
transport of the labeled experiment, the anchor scalar, the curvature
factors that may enter an ensemble, and the failure of Gibbs regularity after
a global label is erased. The label rates found here are the bulk scores of
the Fisher-density statements in \cref{ch:thermodynamic-fisher-draft}.

The finite ensembles, their two normalization modes, and their finite VFE
identities are those of \cref{def:rg-geometry-ensemble,prop:rg-geometry-weights}
in the main manuscript; this chapter adds only what the lattice and the
volume limit require. The lattice notation of
\cref{ch:constructive-spatial-rg} is used throughout, with branch subscripts:
$b_\gamma=\tanh h_\gamma$ and $t_\gamma=\tanh j_\gamma$ are the field and
plaquette strengths of branch $\gamma\in\mathfrak G$, $\tau$ is the decay
exponent of \cref{thm:rg-spatial-constrained-specification}, and $L$ is the
side of the square $\Lambda_L=\{0,\ldots,L-1\}^2$ of cells. For a cell
interaction $\mathscr U$, $\mathsf m_0(\mathscr U)$ and $\mathsf m_1(\mathscr U)$ are its norm
and first diameter moment and $\mathfrak p(\mathscr U)$ its pressure; the specific
entropy, energy, relative entropy, and VFE are $\mathfrak s$,
$\mathfrak e_{\mathscr U}$, $\mathfrak h$, and $\mathfrak f$. Record kernels of the
finite ensemble are written $K_\gamma(o\mid x)$.

The lattice is a realization, and whether its sites are agents at $c_*$ or
points of a base is \cref{open:rg-lattice-typing}. The chapter does not claim
any observation statement for records other than the constant all-positive
record of \cref{hyp:rg-pressure-positive-record}, any conclusion drawn by
conditioning on a null infinite event, or any derivative statement, which
belongs to \cref{ch:thermodynamic-fisher-draft}. The quenched evidence
density of general records is \cref{open:rg-pressure-quenched-record}. Uniform invertibility at
unbounded depth, simultaneous volume and depth limits of the labeled
experiment, self-consistent curvature actions, critical universality, and
continuum field measures are collected in \cref{open:rg-spatial-generality}.
\status{NOT-CLAIMED}

\section{Transport of reference, scalar, and experiment}
\label{sec:rg-geometry-transport}

\theoremheading{Finite transport of the labeled experiment}{thm:rg-geometry-transport}
In the ensemble of \cref{def:rg-geometry-ensemble}, let $C_\gamma(x,y)$ be,
for each label, a normalized channel that is independent of recognition and
of parameters and retains $\gamma$; suppress the label. On targets reachable
from the prior define
\begin{equation}
 \nu'(y)=\sum_x\nu(x)C(x,y),\qquad
 r'(y)=\frac{\sum_x\nu(x)r(x)C(x,y)}{\nu'(y)},\qquad
 p'(y)=\frac{\nu'(y)r'(y)}{Z}.
 \label{eq:rg-geometry-raw-push}
\end{equation}
Then $\sum_y\nu'(y)r'(y)=Z$, so $p'=pC$. The likelihood induced through the
prior reverse row,
\begin{equation}
 \overleftarrow C_p(x\mid y)=\frac{p(x)C(x,y)}{p'(y)},\qquad
 K'(o\mid y)=\sum_x\overleftarrow C_p(x\mid y)K(o\mid x),
 \label{eq:rg-geometry-likelihood}
\end{equation}
preserves the evidence of every record, the selected posterior is pushed by
the same channel, and reference, prior, posterior, recognition, and the
observation joint compose along finite channel words. The reference reverse
row $\nu(x)C(x,y)/\nu'(y)$ induces $K'(\cdot\mid y)$ for every likelihood $K$
exactly when $r$ is constant on the reverse fiber of $y$.
\status{ESTABLISHED}

\noindent\emph{Proof.}
The reference rows sum to one, and strict positivity of $\nu$ and $r$ makes
$\nu'(y)>0$ exactly on the reachable support of $pC$. Multiplying and
summing in \cref{eq:rg-geometry-raw-push} proves the partition identity.
The row $\overleftarrow C_p$ is normalized and
$p'(y)K'(o\mid y)=\sum_xp(x)C(x,y)K(o\mid x)$; summing over $y$ preserves the
evidence, and the posterior statement is the finite case of
\cref{thm:cg-evidence-preserving-channel}. For a following channel $D$,
interchanging the finite sums over $x$ and $y$ gives the composite $CD$. The
ratio of the two reverse rows at $x$ is $r(x)/r'(y)$, which is one on the
reverse fiber exactly when $r$ is constant there; if the rows differ at some
$x_0$, the two-outcome likelihood with $K(o_1\mid x)=\mathbf1_{\{x=x_0\}}$
receives different values from them. On a null target the raw
quotient and the reverse rows are undefined; that column contributes zero,
even if the next channel has a nonzero row there. $\square$

\propositionheading{Anchor scalar cocycle}{prop:rg-geometry-scalar}
Choose a reachable reference anchor $x_*$ and write $r=\mathsf z_0\widehat r$
with $\widehat r(x_*)=1$ and a physical scalar $\mathsf z_0>0$. Denote the raw
push of $\widehat r$ relative to $\nu C$ by $C_\nu\widehat r$, and choose a
reachable target anchor $y_*$. Then
\begin{equation}
 \mathsf z_C=(C_\nu\widehat r)(y_*),\qquad
 \widehat r'=\frac{C_\nu\widehat r}{\mathsf z_C},\qquad
 \mathsf z_0'=\mathsf z_0\mathsf z_C,
 \label{eq:rg-geometry-anchor}
\end{equation}
and, with the same final reference and anchor,
\begin{equation}
 \mathsf z_{CD}(\nu,\widehat r)
 =\mathsf z_C(\nu,\widehat r)
   \mathsf z_D(\nu C,\widehat r').
 \label{eq:rg-geometry-scalar-cocycle}
\end{equation}
The global-factor weight must absorb $\mathsf z_0'$ if only $\widehat r'$ is used
as the branch factor, whereas the scalar cancels within a hierarchical
branch. \status{ESTABLISHED}

\noindent\emph{Proof.}
The numerator of the two-stage raw push is
$\sum_{x,y}\nu(x)\widehat r(x)C(x,y)D(y,z)$ and its denominator is
$(\nu CD)(z)$. The one-stage push relative to $\nu C$ applied to
$\widehat r'$ has the same numerator divided by $\mathsf z_C$, so linearity
pulls the first anchor scalar outside the second push, and evaluation at the
final anchor gives the cocycle. The product $\mathsf z_0'\widehat r'$ recovers
the raw factor, and substitution into
\cref{eq:rg-geometry-hierarchy,eq:rg-geometry-factor} gives the two weight
rules. The identity channel has scalar one only for an unchanged anchored
representative. $\square$

Anchoring changes a representative of the raw factor; it does not change a
principal frame.

\section{Curvature factors and the gauge action}
\label{sec:rg-geometry-curvatures}

\definitionheading{Admitted curvature factor}{def:rg-geometry-curvature-factor}
Gauge curvature is the curvature $F_\omega$ of a principal connection $\omega$,
base curvature that of a nondegenerate metric on $\mathcal C$, and Fisher
curvature that of a regular statistical family with its Fisher metric. A
curvature factor of an ensemble is a function of one of these objects,
together with its pairing, metric contraction, measure, coupling, and a
finite normalizer, entering the ensemble through a raw factor, a label
weight, or a likelihood. For a unitary representation of dimension $n$, a
Riemannian metric on $\mathcal C$ that raises both base indices, its volume
measure, and a coupling $\kappa_{\mathrm g}\geq0$, the gauge action is
\begin{equation}
 \mathscr S_{\mathrm g}=
 -\frac{\kappa_{\mathrm g}}{4n}\int_{\mathcal C}
 \operatorname{Re}\Tr(F_{\mu\nu}F^{\mu\nu})d\operatorname{vol}_{\mathcal C}.
 \label{eq:rg-geometry-gauge-action}
\end{equation}

\remarkheading{Sign of the gauge action}{rem:rg-geometry-gauge-action-sign}
For anti-Hermitian represented curvature and a Riemannian contraction,
$\operatorname{Re}\Tr(F_{\mu\nu}F^{\mu\nu})\leq0$ pointwise, so the integrand
of \cref{eq:rg-geometry-gauge-action} is nonnegative and
$\mathscr S_{\mathrm g}\geq0$. \status{ESTABLISHED}

\noindent\emph{Proof.}
In an orthonormal frame of the base metric, $F^{\mu\nu}=F_{\mu\nu}$, and for
an anti-Hermitian matrix $X$ one has
$-\operatorname{Re}\Tr(XX)=\Tr(X^*X)\geq0$, with equality only at $X=0$.
Summing over $\mu,\nu$ proves the sign. $\square$

A nonnegative action normalizes a law on connections only together with a
reference measure and a finite nonzero partition function. The three
curvatures have different underlying spaces and are not identified here.

\section{Erasing a global label can leave the Gibbs class}
\label{sec:rg-geometry-erasure}

\theoremheading{Product-branch erasure can destroy Gibbs regularity}
{thm:rg-geometry-erasure}
On a countably infinite lattice choose $0<p<q<1$ and $0<w<1$, and let
$\mu_p,\mu_q$ be independent Bernoulli product laws. The mixture
$\mu=w\mu_p+(1-w)\mu_q$ has no singleton conditional version that is
continuous in the product topology at any exterior. In particular it has no
absolutely summable finite-alphabet Gibbs interaction, and both branches can
lie in an arbitrarily small neighborhood of the fair product law. This is
the Bernoulli-product case of the convex-combination theorem of
\citet[Sec.~4.5.1, Prop.~4.12 and Cor.~4.13]{VanEnterFernandezSokal1993},
which gives the absence of a Feller, and for a finite alphabet of a
quasilocal, specification; the discontinuity at every exterior is the
sharper form proved here. \status{ESTABLISHED}

\noindent\emph{Proof.}
Enumerate the exterior sites and let $T_r$ be the event that their empirical
mean tends to $r$. For centered Bernoulli variables
$\E(\sum_{k=1}^n(X_k-r))^4\leq n^2/4$, so a mean deviation above $\epsilon$
has probability at most $1/(4n^2\epsilon^4)$, which is summable; hence
$\mu_r(T_r)=1$, also after removing finitely many sites. The exterior
therefore identifies the branch almost surely, and by independence within a
branch a singleton conditional version equals $p$ on $T_p$ and $q$ on $T_q$.
Every nonempty finite exterior cylinder has positive mass under both
branches, so it meets both $T_p$ and $T_q$ in positive mixture mass, and any
other version agrees off a null set; every cylinder therefore carries both
values, and continuity at any exterior would require a cylinder with
oscillation below $q-p$. An absolutely summable interaction gives a uniformly
convergent singleton energy sum of cylinder functions, whose normalized
exponential would be a continuous version. The one-site fields
$\tfrac12\log[r/(1-r)]$ tend to zero as $r\to1/2$, which gives the last
assertion. $\square$

The labeled joint itself remains normalized, and its finite erasure is an
ordinary stochastic push with the posterior-fiber loss of
\cref{thm:rg-exact-coarse-vfe}. What fails is membership of the erased
mixture in the Gibbs class of each branch.

\section{Carrier, boundary convention, and pressure}
\label{sec:rg-pressure-carrier}

\definitionheading{Native cell model}{def:rg-pressure-native}
A cell $x\in\mathbb Z^2$ carries its outgoing signs $z_{x,1},z_{x,2}$. On
$\Lambda_L$ use the normalized counting reference $4^{-L^2}$, all $2L^2$
outgoing signs, and the $(L-1)^2$ plaquettes whose four edges are outgoing
edges of cells in $\Lambda_L$. Branch $\gamma$ has raw factor
\begin{equation}
 r_{\gamma,L}(z)=\prod_{e\in\Lambda_L\times\{1,2\}}(1+b_\gamma z_e)
 \prod_{\square\subset\Lambda_L}(1+t_\gamma z_{\partial\square}),
 \qquad 0<b_\gamma<1,\quad 0<t_\gamma<1/12,
 \label{eq:rg-pressure-native-factor}
\end{equation}
where $z_{\partial\square}$ is the product of the four signs of the
plaquette, and the finite catalog of strengths lies in a compact subdomain.
An all-touching fixed-exterior model includes every plaquette meeting the
free cells and is a different finite boundary convention. Independent
uniform gauge frames contribute reference mass one. For local limits the
coarse square is centered at $-\lfloor L/2\rfloor$ in both coordinates and
the fine square of side $BL$ at $-B\lfloor L/2\rfloor$.

\remarkheading{Centered local limits}{rem:rg-pressure-centered-limits}
With this centering every fixed cylinder has diverging distance to the
boundary, and the free, fixed-boundary, and infinite-law-marginal
approximations of a branch converge on every fixed cylinder to its unique
Gibbs law. A box with a fixed corner does not exhaust $\mathbb Z^2$ and gives
no such statement. \status{ESTABLISHED}

\noindent\emph{Proof.}
Condition each of the three laws on the signs outside a centered inner square
one interaction layer inside $\Lambda_L$. Every plaquette meeting the inner
square is present, so each conditional is the specification of the inner
square with the conditioning exterior. The estimate of
\cref{thm:rg-spatial-constrained-specification} is uniform over exteriors;
integrating it over the three exterior laws and letting the distance from a
fixed cylinder to the boundary diverge proves convergence to the unique
Gibbs law of \cref{thm:rg-spatial-source-law}. $\square$

\theoremheading{Constructive pressure and surface comparison}
{thm:rg-pressure-boundary}
Let $\mathscr U$ be a translation-invariant cell interaction with
\begin{equation}
 \mathsf m_0(\mathscr U)=\sum_{A\ni0}\|\mathscr U_A\|_\infty<\infty,\qquad
 \mathsf m_1(\mathscr U)=\sum_{A\ni0}\operatorname{diam}(A)\|\mathscr U_A\|_\infty<\infty,
 \label{eq:rg-pressure-moments}
\end{equation}
using the Chebyshev diameter of cells. Write $Z_L$ for its free partition
relative to normalized counting and $\mu_{\Lambda_L}$ for the marginal of a
Gibbs law with that interaction. Then
\begin{align}
 \mathfrak p(\mathscr U)&=\lim_{L\to\infty}L^{-2}\log Z_L,
 &\left|\mathfrak p(\mathscr U)-n^{-2}\log Z_n\right|&\leq\frac{4\mathsf m_1(\mathscr U)}n,
 \label{eq:rg-pressure-limit}\\
 \left|\log\frac{\mu_{\Lambda_L}(z)}{p_L^{\mathrm{free}}(z)}\right|
 &\leq8\mathsf m_1(\mathscr U)L,
 \label{eq:rg-pressure-surface}
\end{align}
and the same mass-ratio bound holds for any fixed exterior. The native
interaction, with singleton energy $-\log[(1+bz_1)(1+bz_2)]$ and plaquette
energy $-\log(1+tz_{\partial\square})$, satisfies
\begin{equation}
 \mathsf m_0\leq\frac{2b}{1-b}+\frac{3t}{1-t},\qquad
 \mathsf m_1\leq\frac{3t}{1-t}.
 \label{eq:rg-pressure-native-moment}
\end{equation}
Existence of the pressure for absolutely summable translation-invariant
interactions is standard \citep[Chs.~15 and~16]{Georgii1988}; the
first-moment hypothesis supplies the explicit rates $4\mathsf m_1/n$ and
$8\mathsf m_1L$, which are proved here without an infinite partition
function. \status{ESTABLISHED}

\noindent\emph{Proof.}
Put $T(d)=\sum_{A\ni0,\operatorname{diam}(A)\geq d}\|\mathscr U_A\|_\infty$. A
support crossing the square and containing $x$ has diameter at least
$\operatorname{dist}(x,\Lambda_L^c)$, and at most $4L$ cells lie at each
positive boundary distance, so the crossing energy is at most
$4L\sum_{d\geq1}T(d)=4L\mathsf m_1$. Adding the exterior changes the free
Hamiltonian by at most $4L\mathsf m_1$, its log partition by the same amount,
and its normalized log mass by at most $8L\mathsf m_1$; integration over an
exterior law preserves the multiplicative bounds. Tile a large square by
$n$-squares and a leftover region $D$. Interactions touching $D$ have total
norm at most $\mathsf m_0|D|$, and crossings of the small squares cost at most
$4\mathsf m_1n$ per tile. Assigning independent normalized counting to $D$
and comparing partition sums exponentially, with $n$ fixed and the leftover
fraction tending to zero, puts both limiting upper and lower pressures within
$4\mathsf m_1/n$ of $n^{-2}\log Z_n$; their difference is at most
$8\mathsf m_1/n$ for every $n$, so they agree. For the native interaction a
cell lies in its singleton support, of diameter zero, and in three plaquette
supports of diameter one. The bounds $-\log(1-u)\leq u/(1-u)$, obtained by
integrating $(1-u)^{-1}$, and $\log(1+u)\leq-\log(1-u)$ for $0\leq u<1$ give
\cref{eq:rg-pressure-native-moment}. $\square$

\section{Pressure of the coarse interaction}
\label{sec:rg-pressure-coarse}

Fix an integer block size $B$ before any volume limit, and let $\mathscr U_B$ be
the canonical vacuum interaction of the pushed law $\mu_B$ of
\cref{def:rg-spatial-vacuum-potential,thm:rg-spatial-generated-interaction},
with vacuum $(+,+)$. Its pressure exists by \cref{thm:rg-pressure-boundary},
because $\mathsf m_1(\mathscr U_B)$ is bounded by \cref{eq:rg-spatial-vacuum-norm}.
The following explicit shell bounds use one positive-coupling domain for all
blocks and all branches, including their observation-shifted fields.

\propositionheading{Explicit coarse shell bounds}{prop:rg-pressure-coarse-shells}
Let $b_+<1$ and $t_+<1/12$ bound the field and plaquette strengths of the
branches and of their shifted fields in \cref{prop:rg-pressure-sensor}. Take
the decay exponent and constants
\begin{align}
 e^{\tau}&=\frac{1+12t_+}{24t_+},&
 \mathrm q&=12t_+e^\tau=\frac{1+12t_+}{2},& \mathrm c_1&=\frac{16\mathrm q}{1-\mathrm q},
 \label{eq:rg-pressure-locality-constants}\\
 \rho&=\frac{b_+(1+t_+^2)+2t_+}{1+t_+^2+2b_+t_+},&
 p_{\min}&=\frac{(1-\rho)^2}{4},&\mathrm c_3&=\frac{2\mathrm c_1}{p_{\min}}.
 \label{eq:rg-pressure-positive-constants}
\end{align}
Then $\tau>0$, $0<\mathrm q<1$, $0<\rho<1$, and for every $B\geq1$ the canonical
interaction satisfies
\begin{equation}
 \|\mathscr U_{B,(x,0)}\|_\infty\leq\frac{4\rho^B}{1-\rho^2},\qquad
 \|\mathscr U_{B,(x,r)}\|_\infty\leq \mathrm c_3B^2re^{-\tau r}\quad(r\geq1),
 \label{eq:rg-pressure-shell}
\end{equation}
where $(x,r)$ indexes the radius-$r$ support anchored at $x$, and
\begin{equation}
 \mathsf m_1(\mathscr U_B)\leq\overline{\mathsf m}_B
 :=2\mathrm c_3B^2\sum_{r\geq1}r^2(2r+1)^2e^{-\tau r}<\infty.
 \label{eq:rg-pressure-coarse-moment}
\end{equation}
\status{ESTABLISHED}

\noindent\emph{Proof.}
Since $t_+<1/12$, the displayed exponent is positive and
$\mathrm q=(1+12t_+)/2<1$; as $\mathrm q$ and $\mathrm c_1$ increase with the plaquette strength, the
modulus \cref{eq:rg-spatial-coarse-modulus} holds with these constants for
every branch, and bounds the target variation by
$\mathrm c_1|A|B^2(R+1)e^{-\tau(R+1)}$ when two exteriors agree within radius $R$.
Conditional on the other fine signs, a path consists of independent signs
with field magnitude at most $\operatorname{atanh}b_++2\operatorname{atanh}t_+$,
so revealing the two target paths successively puts every one-cell atom in
$[(1-\rho^B)^2/4,(1+\rho^B)^2/4]$, hence above $p_{\min}$. The singleton
term of the canonical interaction is a conditional log-odds of two such atoms,
bounded by $4\operatorname{atanh}\rho^B\leq4\rho^B/(1-\rho^2)$. The radius-$r$
term is the difference of two successive log-odds whose exteriors agree
within radius $r-1$; each of the two probabilities moves by at most
$\mathrm c_1B^2re^{-\tau r}$, and the logarithm is $1/p_{\min}$-Lipschitz on
$[p_{\min},1]$, which gives the second bound. At most $(2r+1)^2$
translated radius-$r$ supports contain a fixed cell, each of diameter at most
$2r$; summing proves \cref{eq:rg-pressure-coarse-moment}. $\square$

The explicit bound $\overline{\mathsf m}_B$ grows with $B$, whereas the bound of
\cref{eq:rg-spatial-vacuum-norm} decays; either supplies the first-moment
hypothesis of \cref{thm:rg-pressure-boundary}.

\section{Label weights, selection, and ties}
\label{sec:rg-pressure-selection}

\hypothesisheading{Controlled label growth}{hyp:rg-pressure-label-growth}
For the global-factor ensemble,
$\log a_{\gamma,L}=L^2\alpha_\gamma+\delta_{\gamma,L}$ with
$|\delta_{\gamma,L}|\leq U_\gamma(L)=o(L^2)$ and a supplied modulus for this
$o(L^2)$ bound, and each branch uses its matching free partition and law, or
its matching fixed-exterior pair. This excludes volume-dependent weights
without a controlled subextensive part. It is the value-level lattice form of
the growth hypothesis \cref{eq:rg-thermo-label-effective-growth} of the
finite-mixture Fisher asymptotics in \cref{sec:cg-retained-label-asymptotics};
the derivative statements of \cref{ch:thermodynamic-fisher-draft} impose that
$C^3$ form instead.

\theoremheading{Strict selection and the exact tie condition}
{thm:rg-pressure-selection}
Assume \cref{hyp:rg-pressure-label-growth}. Write
$\mathsf s_\gamma=\alpha_\gamma+\mathfrak p_\gamma$ and
$W_{\gamma,L}=a_{\gamma,L}Z_{\gamma,L}/\sum_{\gamma'}a_{\gamma',L}Z_{\gamma',L}$,
and let $\overline{\mathsf m}_\gamma$ be a finite bound for the first moment
$\mathsf m_1(\mathscr U_\gamma)$ of the interaction of branch $\gamma$, for example
$3t_\gamma/(1-t_\gamma)$ from \cref{eq:rg-pressure-native-moment} for a
native branch or $\overline{\mathsf m}_B$ of
\cref{eq:rg-pressure-coarse-moment} for a canonical coarse branch. For free
boundaries put $\Delta_\gamma(L)=U_\gamma(L)+4\overline{\mathsf m}_\gamma L$,
and replace $4$ by $8$ for fixed exteriors. If $\gamma_*$ uniquely maximizes
$\mathsf s_\gamma$, then
\begin{equation}
 1-W_{\gamma_*,L}\leq\sum_{\gamma\ne\gamma_*}
 \exp\left[-L^2(\mathsf s_{\gamma_*}-\mathsf s_\gamma)
                  +\Delta_\gamma(L)+\Delta_{\gamma_*}(L)\right]\longrightarrow0.
 \label{eq:rg-pressure-selection-bound}
\end{equation}
If instead $S$ is the full maximal-score set and
$\widehat\delta_{\gamma,L}=\delta_{\gamma,L}+\log Z_{\gamma,L}-L^2\mathfrak p_\gamma$,
the labeled local law converges exactly when the finite vector
\begin{equation}
 \left(\frac{e^{\widehat\delta_{\gamma,L}}}
 {\sum_{\gamma'\in S}e^{\widehat\delta_{\gamma',L}}}\right)_{\gamma\in S}
 \label{eq:rg-pressure-tie-vector}
\end{equation}
converges; pressure equality alone does not determine it.
\status{ESTABLISHED}

\noindent\emph{Proof.}
The pressure error bounds give
$|\log(a_{\gamma,L}Z_{\gamma,L})-L^2\mathsf s_\gamma|\leq\Delta_\gamma(L)$.
Dividing each competing weight by that of $\gamma_*$ and bounding the common
denominator below by its largest term proves
\cref{eq:rg-pressure-selection-bound}. For a tied maximal set the complement,
if nonempty, has a positive gap and hence vanishing mass, and within $S$ the
relative weights are exactly \cref{eq:rg-pressure-tie-vector}. Branch
cylinder limits exist by \cref{rem:rg-pressure-centered-limits}. Finite
summation proves sufficiency, and testing the cylinder consisting only of the
label proves necessity. $\square$

The Fisher-tensor counterparts of the strict and tie statements are the
lattice-free statements of \cref{sec:cg-retained-label-asymptotics}; their
lattice instances are collected in \cref{ch:thermodynamic-fisher-draft}.

\remarkheading{Ties are not fixed by pressures}{rem:rg-pressure-ties}
For two identical positive-coupling branches with $a_{1,L}=v_L/Z_L$ and
$a_{2,L}=1/Z_L$, the choice of $v_L$ alternating between $1$ and $2$ gives
identical bulk scores and label weights alternating between $1/2$ and $2/3$,
and $v_L=e^{-L}$ removes one label with zero bulk gap. A fixed hierarchy
$w_\gamma\mu_\gamma$ is a finite sum of probability laws, reproduced at each
finite volume by the factor weights $a_{\gamma,L}=w_\gamma/Z_{\gamma,L}$ of
\cref{prop:rg-geometry-weights}, but compensation by
$e^{-L^2\mathfrak p_\gamma}$ alone leaves the surface term
$\log Z_{\gamma,L}-L^2\mathfrak p_\gamma$: multiplying one of two identical
raw branches by $e^L$ leaves both pressures unchanged and makes this term
select that label. A bulk counterterm therefore does not restore prescribed
tied weights in general. \status{ESTABLISHED}

\noindent\emph{Proof.}
The weights are proportional to $a_{\gamma,L}Z_{\gamma,L}$, which gives
$v_L/(1+v_L)\in\{1/2,2/3\}$ and $e^{-L}/(1+e^{-L})\to0$, while
$L^{-2}\log v_L\to0$ leaves both bulk scores equal. The branch $e^L$ times a
raw factor has partition $e^LZ_{\gamma,L}$ and the same pressure, and its
relative weight $e^L\to\infty$ selects it. The example concerns an explicit
scalar factor, not the sign of a native boundary correction. $\square$

\section{Observation, entropy, and specific VFE}
\label{sec:rg-pressure-vfe}

\hypothesisheading{Constant all-positive observation record}
{hyp:rg-pressure-positive-record}
Each branch has the independent edge sensor
$K^{\mathrm{sens}}_\gamma(o\mid z)=(1+\eta_\gamma oz)/2$ with
$0<\eta_\gamma<1$, and the observation conditioned on is the single finite
record in which all $2L^2$ outgoing edge sensors of $\Lambda_L$ report
$o=+$. This record is a choice, and it excludes generic and typical records.
Under every branch it has probability at most
$[(1+\eta_+)/2]^{2L^2}$, with $\eta_+=\max_\gamma\eta_\gamma$. A negative entry
$o_e=-$ multiplies the prior factor by
$(1+bz_e)(1-\eta z_e)/2=A_-(1+b'_-z_e)$ with $A_-=(1-b\eta)/2$ and
$b'_-=(b-\eta)/(1-b\eta)$, which is negative whenever $b<\eta$ and then
leaves the positive-field domain of
\cref{def:rg-pressure-native,thm:rg-spatial-constrained-specification}; a
general record produces a site-dependent field of mixed sign, whose evidence
density is the subject of \cref{open:rg-pressure-quenched-record}. The
hypothesis is used in \cref{prop:rg-pressure-sensor} and in
\cref{thm:rg-pressure-vfe-density,prop:rg-pressure-vfe-minimum,thm:rg-pressure-induced-vfe}.

\propositionheading{Positive sensor and posterior branch}
{prop:rg-pressure-sensor}
Under \cref{hyp:rg-pressure-positive-record}, define
\begin{equation}
 A_\gamma=\frac{1+b_\gamma\eta_\gamma}{2},\qquad
 b'_\gamma=\frac{b_\gamma+\eta_\gamma}{1+b_\gamma\eta_\gamma},\qquad
 Z^{\mathrm{obs}}_{\gamma,L}
 =A_\gamma^{2L^2}Z_L(b'_\gamma,t_\gamma).
 \label{eq:rg-pressure-sensor}
\end{equation}
The selected free or fixed-boundary posterior is the matching shifted branch,
and reweighting infinite-law marginals gives the same centered local
posterior limit. The evidence density of branch $\gamma$ is
\begin{equation}
 e_\gamma=2\log A_\gamma+
 \mathfrak p(b'_\gamma,t_\gamma)-\mathfrak p(b_\gamma,t_\gamma).
 \label{eq:rg-pressure-evidence}
\end{equation}
For fixed positive hierarchical weights the total evidence density is
$\bar e=\max_\gamma e_\gamma$ and the posterior label scores are
$\mathsf s^o_\gamma=e_\gamma$. For global weights under
\cref{hyp:rg-pressure-label-growth}, with
$\mathfrak p^{\mathrm{obs}}_\gamma=\mathfrak p_\gamma+e_\gamma$, they are
\begin{equation}
 \bar e=\max_\gamma(\alpha_\gamma+\mathfrak p_\gamma^{\mathrm{obs}})
   -\max_\gamma(\alpha_\gamma+\mathfrak p_\gamma),\qquad
 \mathsf s^o_\gamma=\alpha_\gamma+\mathfrak p_\gamma^{\mathrm{obs}}.
 \label{eq:rg-pressure-evidence-global}
\end{equation}
These are limits of finite observations, not conditioning on the null
infinite all-positive event. \status{ESTABLISHED}

\noindent\emph{Proof.}
The identity $(1+bz)(1+\eta z)/2=A(1+b'z)$ follows from $z^2=1$; multiplying
it over the observed free edges leaves the plaquette factors unchanged. The
shifted field remains below one and the plaquette strength below $1/12$. For
an infinite-law marginal, the finite reweighted full-space law conditioned on
the exterior of the window is the shifted all-touching specification, and
integrating the uniform interior comparison over that reweighted exterior
proves local convergence; a surface log-ratio bound alone would not. For
densities, the prior log-ratio between marginal and free law is at most
$8\mathsf m_1L$; multiplying by the same positive likelihood bounds the
evidence log-ratio by $8\mathsf m_1L$ and the posterior log-ratio by
$16\mathsf m_1L$, and the comparison for the shifted interaction, of first
moment $\mathsf m_1'$, adds $8\mathsf m_1'L$. These
vanish after division by $L^2$ and give \cref{eq:rg-pressure-evidence}. The
unnormalized posterior label masses are $w_\gamma E_{\gamma,L}$ for the
hierarchy and $a_{\gamma,L}Z^{\mathrm{obs}}_{\gamma,L}$ for global factors,
and the logarithm of a finite exponential sum differs from its largest term
by at most $\log|\mathfrak G|$, which gives the label rates. The infinite
all-positive event is null because the finite likelihoods are bounded by
$[(1+\eta_+)/2]^{2L^2}\to0$. $\square$

Posterior ties require their own subextensive ratios, as in
\cref{eq:rg-pressure-tie-vector}.

\openproblemheading{Quenched evidence density for general records}
{open:rg-pressure-quenched-record}
Let the record $o=(o_e)$ on all outgoing edges be drawn from a stationary
ergodic law, for example from the model itself, and condition on its
restriction to $\Lambda_L$. The problem is whether $L^{-2}\log E_L(o)$, the
posterior label rates, and the specific VFE of
\cref{thm:rg-pressure-vfe-density} converge almost surely to
record-independent quenched limits, and whether
\cref{prop:rg-pressure-vfe-minimum} holds for those limits. Settling it
requires an every-exterior comparison and a pressure identity for the
site-dependent shifted fields in $(-1,1)$, of both signs, that a general
record produces, together with an ergodic theorem for the record-dependent
evidence. The obstruction is that
\cref{thm:rg-spatial-constrained-specification} and every construction of
this chapter are stated for one constant field $0<b<1$ and use translation
invariance of the branch law, which a fixed realization of a random record
breaks.

\theoremheading{Stationary specific relative entropy}
{thm:rg-pressure-specific-kl}
For a stationary law $Q$ on the four-state cell lattice with square entropy
$\operatorname{Ent}(Q_{\Lambda_L})$, using $0\log0=0$, and a Gibbs law $\mu$
whose interaction satisfies \cref{eq:rg-pressure-moments},
\begin{align}
 \mathfrak s(Q)&=\lim_{L\to\infty}L^{-2}\operatorname{Ent}(Q_{\Lambda_L})
       =\inf_{n\geq1}n^{-2}\operatorname{Ent}(Q_{\Lambda_n}),
 \label{eq:rg-pressure-entropy}\\
 \mathfrak e_{\mathscr U}(Q)&=\sum_{A\ni0}\frac{\E_Q\mathscr U_A}{|A|},
 \label{eq:rg-pressure-energy}\\
 \mathfrak h(Q\Vert\mu)&:=\lim_{L\to\infty}L^{-2}
       \KL(Q_{\Lambda_L}\Vert\mu_{\Lambda_L})
       =\log4-\mathfrak s(Q)+\mathfrak e_{\mathscr U}(Q)+\mathfrak p(\mathscr U).
 \label{eq:rg-pressure-specific-kl}
\end{align}
The limits need neither ergodicity nor positive recognition atoms, and they
concern specific relative entropy, not relative entropy between infinite
laws. The formula is the finite-alphabet form of the specific entropy,
specific energy, and specific free energy of
\citet[Secs.~15.2 and~15.3, Thm.~15.30]{Georgii1988}; the first-moment
hypothesis supplies the explicit surface rate used in the proof.
\status{ESTABLISHED}

\noindent\emph{Proof.}
Tile a large square by $n$-squares. Subadditivity of entropy and
stationarity bound its entropy by the number of tiles times
$\operatorname{Ent}(Q_{\Lambda_n})$ plus $|D|\log4$, so the limiting upper density is at most
every $n$-square density, whereas their infimum is at most the limiting
lower density. Allocating each interaction equally among its sites,
absolute summability and stationarity give \cref{eq:rg-pressure-energy},
with omitted boundary allocations bounded by $4\mathsf m_1L$. The free
relative entropy equals
$-\operatorname{Ent}(Q_{\Lambda_L})+L^2\log4+\E_Q\mathscr H_L+\log Z_L$, where
$\mathscr H_L$ is the free Hamiltonian, and changing
the comparison to the Gibbs marginal costs at most $8\mathsf m_1L$. Dividing
by $L^2$ and applying \cref{thm:rg-pressure-boundary} proves the formula.
$\square$

\remarkheading{No uniform finite-query entropy bound}
{rem:rg-pressure-finite-query}
Stationarity alone admits no procedure that, after finitely many cylinder
queries, returns the specific entropy with error below $\log4/2$ for every
stationary law. \status{ESTABLISHED}

\noindent\emph{Proof.}
A uniformly random $K$-square four-symbol array repeated periodically is
stationary, has specific entropy zero, and has the fair independent marginal
on every square of side at most $K$. A terminating procedure receives the same
answers, and so returns the same value, for it and for the fair independent
field once $K$ exceeds every queried coordinate span, whereas their specific
entropies differ by $\log4$; one value cannot lie within $\log4/2$ of both.
$\square$

\theoremheading{Label-retaining specific VFE}
{thm:rg-pressure-vfe-density}
Fix a recognition law $\mathbb Q=q_\gamma Q_\gamma$ with each $Q_\gamma$
stationary, and put $\mathfrak h_\gamma=\mathfrak h(Q_\gamma\Vert\mu_\gamma)$
and $\mathfrak l^Q_\gamma=\E_{Q_\gamma}\sum_{a=1}^2
\log K^{\mathrm{sens}}_\gamma(+\mid z_{0,a})$. With
$\mathsf s_*=\max_\gamma\mathsf s_\gamma$, the specific VFE of the two
ensembles are
\begin{equation}
 \mathfrak f_{\mathrm h}=\sum_\gamma q_\gamma(\mathfrak h_\gamma-\mathfrak l_\gamma^Q),\qquad
 \mathfrak f_{\mathrm f}=\sum_\gamma q_\gamma
       (\mathfrak h_\gamma-\mathfrak l_\gamma^Q+\mathsf s_*-\mathsf s_\gamma).
 \label{eq:rg-pressure-vfe-density}
\end{equation}
They exist under \cref{hyp:rg-pressure-positive-record} and, for global
factors, \cref{hyp:rg-pressure-label-growth}. \status{ESTABLISHED}

\noindent\emph{Proof.}
Disintegrate the finite relative entropy over the label. For the hierarchy
$\KL(q\Vert w)/L^2\to0$, and the branch limits follow from
\cref{thm:rg-pressure-specific-kl}; stationarity evaluates the likelihood sum
per cell exactly. For global factors
$-L^{-2}\log W_{\gamma,L}\to\mathsf s_*-\mathsf s_\gamma$ by the finite
exponential bounds, and the recognition label entropy again vanishes per cell.
Summing finitely many branch terms proves the formulas; matching
fixed-exterior or compared infinite-law-marginal versions have the same
density by the surface bounds. $\square$

\propositionheading{Minimum of the specific VFE}
{prop:rg-pressure-vfe-minimum}
Assume the hypotheses of \cref{thm:rg-pressure-vfe-density}. Let
$\mu^{\mathrm{post}}_\gamma$ be the Gibbs law of the shifted branch
$(b'_\gamma,t_\gamma)$ of \cref{prop:rg-pressure-sensor}, unique and
translation invariant by \cref{thm:rg-spatial-source-law} because
$0<b'_\gamma<1$ and $t_\gamma<1/12$, and put
$\mathfrak h^{\mathrm{post}}_\gamma=\mathfrak h(Q_\gamma\Vert\mu^{\mathrm{post}}_\gamma)$.
With the posterior label scores $\mathsf s^o_\gamma$ and total evidence
density $\bar e$ of \cref{prop:rg-pressure-sensor} and
$\mathsf s^o_*=\max_\gamma\mathsf s^o_\gamma$, every recognition
$\mathbb Q=q_\gamma Q_\gamma$ with stationary branches satisfies, for
$\mathfrak f=\mathfrak f_{\mathrm h}$ or $\mathfrak f=\mathfrak f_{\mathrm f}$,
\begin{equation}
 \mathfrak f=-\bar e+\sum_\gamma q_\gamma
 \left[(\mathsf s^o_*-\mathsf s^o_\gamma)+\mathfrak h^{\mathrm{post}}_\gamma\right],
 \qquad
 \mathfrak h^{\mathrm{post}}_\gamma=\mathfrak h_\gamma-\mathfrak l^Q_\gamma+e_\gamma.
 \label{eq:rg-pressure-vfe-decomposition}
\end{equation}
Consequently
\begin{equation}
 \min_{\mathbb Q}\mathfrak f=-\bar e,
 \label{eq:rg-pressure-vfe-minimum}
\end{equation}
and the minimizers are exactly the recognitions whose label law is supported
on the maximal posterior-score set
$\{\gamma:\mathsf s^o_\gamma=\mathsf s^o_*\}$ and whose charged branches are
$Q_\gamma=\mu^{\mathrm{post}}_\gamma$. When several labels attain
$\mathsf s^o_*$, every label law on them is a minimizer; the density does not
determine limits of the finite posterior label weights, which require the
posterior analogue of \cref{eq:rg-pressure-tie-vector}. At each finite volume
the unique minimizer of the finite VFE is the finite posterior, whose branch
conditionals converge locally to $\mu^{\mathrm{post}}_\gamma$ by
\cref{prop:rg-pressure-sensor}. \status{ESTABLISHED}

\noindent\emph{Proof.}
Fix a boundary convention of \cref{thm:rg-pressure-vfe-density} and let
$\mathbb Q_L$ and $Q_{\gamma,L}$ be the $\Lambda_L$ marginals. Write
$E_{\gamma,L}$ for the branch evidence of the record, $E_L$ for the total
evidence, $W^o_L$ for the posterior label law, $\Pi_{\gamma,L}$ for the
posterior branch, and $F_L(\mathbb Q_L)$ for the finite VFE. The identity
$F_L(\mathbb Q_L)=-\log E_L+\KL(\mathbb Q_L\Vert\boldsymbol\Pi_L)$ of
\cref{thm:elbo-extended-gap} and the chain rule of relative entropy over the
label give
\begin{equation}
 F_L(\mathbb Q_L)=-\log E_L+\KL(q\Vert W^o_L)
 +\sum_\gamma q_\gamma\KL(Q_{\gamma,L}\Vert\Pi_{\gamma,L}).
 \label{eq:rg-pressure-finite-vfe-chain}
\end{equation}
Because the prior branch law $p_{\gamma,L}$ and the record likelihood
$K_L$ are strictly positive,
\begin{equation}
 \KL(Q_{\gamma,L}\Vert\Pi_{\gamma,L})=\KL(Q_{\gamma,L}\Vert p_{\gamma,L})
 -\E_{Q_\gamma}\log K_L+\log E_{\gamma,L}.
 \label{eq:rg-pressure-finite-branch-posterior}
\end{equation}
Divide by $L^2$. The first term on the right tends to $\mathfrak h_\gamma$ by
\cref{thm:rg-pressure-specific-kl}, stationarity makes the second exactly
$-L^2\mathfrak l^Q_\gamma$, and $L^{-2}\log E_{\gamma,L}\to e_\gamma$ by
\cref{prop:rg-pressure-sensor}. For free boundaries $\Pi_{\gamma,L}$ is the
free law of the shifted branch, and for the other conventions its log mass
differs from that law by $O(L)$; hence \cref{thm:rg-pressure-specific-kl},
applied to $\mu^{\mathrm{post}}_\gamma$, identifies the limit of the left side
with $\mathfrak h^{\mathrm{post}}_\gamma$, which proves the second identity in
\cref{eq:rg-pressure-vfe-decomposition}. By
\cref{eq:rg-pressure-sensor,thm:rg-pressure-boundary} and, for global factors,
\cref{hyp:rg-pressure-label-growth}, the logarithms of the unnormalized
posterior label masses are $L^2\mathsf s^o_\gamma+o(L^2)$, so
$L^{-2}\log E_L\to\bar e$ and
$L^{-2}\KL(q\Vert W^o_L)\to\sum_\gamma q_\gamma(\mathsf s^o_*-\mathsf s^o_\gamma)$.
Since $L^{-2}F_L(\mathbb Q_L)\to\mathfrak f$ by
\cref{thm:rg-pressure-vfe-density}, dividing
\cref{eq:rg-pressure-finite-vfe-chain} by $L^2$ proves the first identity.

Every bracket in \cref{eq:rg-pressure-vfe-decomposition} is nonnegative,
because $\mathsf s^o_*\geq\mathsf s^o_\gamma$ and specific relative entropy is
a limit of nonnegative quantities. Hence $\mathfrak f\geq-\bar e$, with equality
exactly when $q$ charges only maximal posterior scores and
$\mathfrak h^{\mathrm{post}}_\gamma=0$ on every charged label. The law
$\mu^{\mathrm{post}}_\gamma$ attains
$\mathfrak h(\mu^{\mathrm{post}}_\gamma\Vert\mu^{\mathrm{post}}_\gamma)=0$.
Conversely, a stationary law whose specific relative entropy with respect to
a translation-invariant Gibbs law of a finite-range specification vanishes
along squares is a Gibbs law of the same specification
\citep[Thm.~15.37]{Georgii1988}, and uniqueness gives
$Q_\gamma=\mu^{\mathrm{post}}_\gamma$. At finite volume
$\KL(\mathbb Q_L\Vert\boldsymbol\Pi_L)\geq0$, with equality only at
$\boldsymbol\Pi_L$, whose branch conditionals are the posterior branches of
\cref{prop:rg-pressure-sensor}. $\square$

\remarkheading{Density convergence and the limiting joint}
{rem:rg-pressure-limit-joint}
If $W_{\gamma,L}\to0$ while $q_\gamma>0$, the limiting labeled joint gives
zero mass to a recognition-positive label, so the relative entropy to that
limiting joint is infinite even on a finite state window, whereas
\cref{eq:rg-pressure-vfe-density} can remain finite; the $e^{-L}$ example of
\cref{rem:rg-pressure-ties} has zero label-rate penalty. Convergence of the
density therefore neither identifies the VFE against a limiting joint nor
requires convergence of the label weights. \status{ESTABLISHED}

\noindent\emph{Proof.}
A recognition-positive atom of zero reference mass makes the relative entropy
infinite. In the $e^{-L}$ example both labels have score $\mathsf s_*$, so
the penalty $\mathsf s_*-\mathsf s_\gamma$ in
\cref{eq:rg-pressure-vfe-density} vanishes, and the alternating example of
\cref{rem:rg-pressure-ties} has a finite density with nonconvergent
weights. $\square$

Erasing the label instead remains subject to \cref{thm:rg-geometry-erasure}.

\section{Scalar composition and the induced VFE density}
\label{sec:rg-pressure-scalar}

\theoremheading{Normalization density along fixed blocking words}
{thm:rg-pressure-scalar}
Let $\mathscr U$ and $\mathscr R^{\mathrm{vac}}_a\mathscr U$ be the canonical vacuum representatives of a
law and of its $a$-path push, each with finite first moment, and let
$\mathsf z_{a,L}$ be the vacuum anchor scalar of the finite raw push. For a
word fixed before the volume limit, with coherent references, anchors, and
centered embeddings,
\begin{align}
 \mathfrak n_a(\mathscr U)&:=\lim_{L\to\infty}L^{-2}\log\mathsf z_{a,L}
        =a^2\mathfrak p(\mathscr U)-\mathfrak p(\mathscr R^{\mathrm{vac}}_a\mathscr U),
 \label{eq:rg-pressure-scalar-density}\\
 \mathfrak n_{ad}(\mathscr U)&=d^2\mathfrak n_a(\mathscr U)
                       +\mathfrak n_d(\mathscr R^{\mathrm{vac}}_a\mathscr U).
 \label{eq:rg-pressure-scalar-word}
\end{align}
\status{ESTABLISHED}

\noindent\emph{Proof.}
Disjoint parity paths push normalized counting to normalized counting, so
$Z_{\mathrm{fine},aL}=\mathsf z_{a,L}\widehat Z_L$ exactly. The fine
comparison between free law and marginal costs $8\mathsf m_1aL$ in log mass,
the push preserves the multiplicative bound, and the coarse comparison adds
$8\mathsf m_1'L$; thus the pushed free law and the canonical coarse free law
differ by at most $8(a\mathsf m_1+\mathsf m_1')L$ in log mass. Both anchored
factors equal one at the vacuum, so evaluating there compares
$\log\widehat Z_L$ with the canonical $\log Z'_L$ by the same amount.
Dividing the partition identity by $L^2$ and using the pressure limits proves
\cref{eq:rg-pressure-scalar-density}; substitution proves
\cref{eq:rg-pressure-scalar-word}, where $d^2$ counts intermediate cells per
final cell, and an unchanged canonical identity has $\mathfrak n_1=0$. $\square$

\definitionheading{Coarse label rates}{def:rg-pressure-coarse-rates}
For a fixed blocking word of area factor $B^2$ and a native label rate
$\alpha$, the native scalar density and the coarse label rate are
\begin{equation}
 \widetilde{\mathfrak n}_B=B^2\mathfrak p_{\mathrm{native}}
                            -\mathfrak p(\mathscr U_B),\qquad
 \alpha'_B=B^2\alpha+\widetilde{\mathfrak n}_B.
 \label{eq:rg-pressure-native-scalar}
\end{equation}

\remarkheading{Native conversion and coarse label scores}
{rem:rg-pressure-native-conversion}
The native raw factor has vacuum logarithm per cell
$2\log(1+b)+\log(1+t)$, and
$\mathfrak p(\mathscr U_1)=\mathfrak p_{\mathrm{native}}-2\log(1+b)-\log(1+t)$; at
$t=0$ the native pressure is zero and the canonical coarse pressure is
$-2\log(1+b^B)$. By \cref{def:rg-pressure-coarse-rates},
$\alpha'_B+\mathfrak p(\mathscr U_B)=B^2(\alpha+\mathfrak p_{\mathrm{native}})$
identically: the coarse label score is the fine score times the area factor,
which is the rate per coarse cell of the labeled push, since that push leaves
every label weight unchanged at each finite volume and a coarse cell carries
$B^2$ fine cells. \status{ESTABLISHED}

\noindent\emph{Proof.}
Dividing the native raw factor by its vacuum value gives a finite-range
vacuum representative whose free partition differs from that of the canonical
representative $\mathscr U_1$ only by a surface term, as both compare to their
common infinite-law marginal at the vacuum; dividing by $L^2$ gives the
conversion. At $t=0$ the raw partition is one, and the pushed coarse cell law
is a product of two independent signs of mean $b^B$, whose canonical vacuum
representative has free partition $(1+b^B)^{-2L^2}$. The last assertion is
algebra. Equality of infinite laws alone would specify neither a raw
normalizer nor a representative conversion. $\square$

\theoremheading{Specific VFE transport of the labeled experiment}
{thm:rg-pressure-induced-vfe}
Retain the label and the fine observation alphabet through a fixed word of
product block size $B$. Let $Q_{\gamma,B}$ be the stationary path push of
$Q_\gamma$ and $\mathfrak h^{\mathrm{post}}_{\gamma,B}$ its specific relative
entropy to the pushed shifted branch. Then
\begin{equation}
 B^2\mathfrak f_{\mathrm{fine}}
   =\mathfrak f_{\mathrm{coarse,induced}}+\mathfrak d_{\mathrm{fib}},\qquad
 \mathfrak d_{\mathrm{fib}}=\sum_\gamma q_\gamma
 \left[B^2\mathfrak h^{\mathrm{post}}_{\gamma,1}
             -\mathfrak h^{\mathrm{post}}_{\gamma,B}\right]\geq0,
 \label{eq:rg-pressure-induced-vfe}
\end{equation}
for each of the two ensembles, using limits of finite quantities rather than
differences of infinite-law relative entropies. \status{ESTABLISHED}

\noindent\emph{Proof.}
On a finite window the induced likelihood is
$K'_\gamma(o\mid y)=E_{\gamma,L}(o)
 p_{\gamma,L}^{\mathrm{post,push}}(y)/p_{\gamma,L}^{\mathrm{prior,push}}(y)$ on
its positive support, by \cref{thm:rg-geometry-transport}, and the finite VFE
difference is the posterior-fiber loss of \cref{thm:rg-exact-coarse-vfe},
which is nonnegative. The coarse recognition marginal is the push of its fine
marginal on the square of side $BL$. The comparison between posterior and
marginal in \cref{prop:rg-pressure-sensor} is $O(BL)$ uniformly in the state
and is preserved by the push, so after division by $L^2$
\cref{thm:rg-pressure-specific-kl} gives both posterior limits. For the
hierarchy the coarse density is
$\sum_\gamma q_\gamma(\mathfrak h^{\mathrm{post}}_{\gamma,B}-B^2e_\gamma)$, and
for global factors it is
$\sum_\gamma q_\gamma[\mathfrak h^{\mathrm{post}}_{\gamma,B}
 +B^2(\mathsf s_*-\alpha_\gamma-\mathfrak p_\gamma^{\mathrm{obs}})]$.
Subtracting these separately finite limits from $B^2$ times the fine formulas
of \cref{prop:rg-pressure-vfe-minimum} cancels the evidence and label rates
and leaves the fiber term, a limit of nonnegative finite losses. $\square$

The identity concerns the likelihood induced from the fine record; a new
local sensor on coarse edges defines a different experiment, to which it does
not apply.

\remarkheading{Label actions and bulk rates}{rem:rg-pressure-experiment-changes}
A prescribed order-one label action has zero bulk rate but can change tied
weights, and an
extensive factor $e^{-L^2\mathscr A_\gamma}$ shifts $\alpha_\gamma$ by
$-\mathscr A_\gamma$. A state-dependent curvature action changes the branch
specification and requires its own admission proof. \status{ESTABLISHED}

\noindent\emph{Proof.}
The bulk rate is $\lim L^{-2}\log a_{\gamma,L}$, which an order-one factor
leaves unchanged and an extensive factor shifts by $-\mathscr A_\gamma$; an
order-one factor multiplies the corresponding entry of
\cref{eq:rg-pressure-tie-vector} and changes the tied weights unless it is
common to all tied labels. A state-dependent factor multiplies the raw
factor, not the weight. $\square$
